
Druhý případ užití je specifikován v tabulce č. \ref{tab:usecase:UC02}.

\begin{UseCaseTable}
  {Přihlášení do systému}
  {UC02}
  {F02}
  {Zahájení přihlášené relace umožňující správu evidovaných subjektů zájmu}
  {Provozovatel monitorovaného prostoru}
  {}
  {Uživatel je v systému evidován v aktivním stavu, není přihlášen}
  {Uživatel je přihlášen}

\UseCaseScenario{Základní scénář}
\UseCaseStep{uživatel}{Otevře přihlašovací stránku.}
\UseCaseStep{systém}{Vrátí přihlašovací formulář.}
\UseCaseStep{uživatel}{Vyplní email a heslo použité při registraci a formulář odešle.}
\UseCaseStep{systém}{Ověří, že zadané údaje odpovídají evidovanému uživateli v aktivním stavu.}
\UseCaseStep{systém}{Zahájí přihlášenou relaci a přesměruje uživatele do administračního rozhraní.}

\UseCaseScenario{Alternativní scénář -- neplatné údaje}
\UseCaseStepCustom{A1}{uživatel, systém}{shodné s kroky 1--3 základního scénáře.}
\UseCaseStepCustom{A2}{systém}{Zjistí, že zadané údaje neodpovídají žádnému evidovanému uživateli, a vyzve uživatele k jejich opětovnému zadání.}
\UseCaseStepCustom{A3}{uživatel}{Zadá údaje znovu.}
\UseCaseStepCustom{A4}{systém}{Pokračuje krokem 4 základního scénáře.}

\UseCaseScenario{Alternativní scénář -- neaktivní účet}
\UseCaseStepCustom{B1}{uživatel, systém}{shodné s kroky 1--3 základního scénáře.}
\UseCaseStepCustom{B2}{systém}{Zjistí, že účet je evidován v neaktivním stavu, a vyzve uživatele k aktivaci účtu přes odkaz v potvrzovacím emailu.}

\end{UseCaseTable}
